% 第8章 线性动态电路暂态过程的时域分析
\chapter{线性动态电路暂态过程的时域分析}

\section{动态电路的暂态过程}
电容元件

在前面

动态

所以
\[
\]
式中    $i$

消去，得到
\[
\]
式是

一阶

综上所述

电路

\section{电路量的初始值}
1.电容电压$u_c$和电感电流$i_L$初始值的确定
设在线性电容，则有
\[
\]
电容电荷的初始值可以表示为
\[
\]
式中等号右端第一项积分表示$t=0_-$时的电荷$q(0_-)$，故
\[
\]
若在，于是得到
\[
\]
式中，称为电容电荷和电压的原始值。

综上所述，若换路瞬间（）电容电流，在$t=0$时是连续变化的或称渐变。

应用对偶原理，的关系
\[
\]
式成立的条件是，

式及其成立的条件所表示的规律称为换路定律。

上述，即
\[
\]
电路换路将引起电容和电感，吸收或释放的功率为
\[
\]
在

2.除电流初始值的确定
瞬间有
 KCL
 KVL
 电阻元件
 电容元件
 电感元件
式中

\section{一阶电路的零输入响应}
可用一阶

在电路的稳态过程中，

1.$\mathrm{RC}$电路的零输入响应
图所示电路

根据微分方程
\[
\]
根据换路定律得初始条件    $u_c(0_+)=u_c(0_-)=U_0$

方程。其求解步骤为

列出特征方程

特征根

微分方程的通解

式中来确定
\[
\]
将上式的解为
\[
\]
再由电阻
\[
\]
由于

将会产生很大的放电电流。

由。现令
\[
\]
称为$\mathrm{RC}$电路的时间常数

下面。

2.$\mathrm{RL}$电路的零输入响应
图所示电路开关原是断开的。其$\mathrm{KVL}$方程为
\[
\]
根据换路定律，得
\[
\]
式，其特征方程和特征根分别为
\[
\]
\[
\]
微分方程的通解是
\[
\]
式中$=\frac{L}{R}$

称为$\mathrm{RL}$电路的时间常数

积分常数$A$
\[
\]
将，得
\[
\]
根据电阻或电感的元件方程求得
\[
\]
式表明

暂态过程

\section{阶跃函数和冲激函数}
1.单位阶跃函数
考察

单位阶跃函数定义为
\[
\]
无量纲

如果，即
\[
\]
的波形

用，记为
\[
\]
$G(t)$的波形

$G(t)$的幅值为。例如
\[
\]
\[
\]
分别表示

2.单位冲激函数
冲激函数，其定义为
\[
\]
$t=0$时，

单位冲激函数

例如

短路

在图

冲激函数

\section{一阶电路的阶跃响应}
若电路

设方程为
\[
\]
将,得
\[
\]
$u_c$的初始值
\[
\]
方程，即
\[
\]
（1）

图，即
\[
\]
（2）

与方程为
\[
\]
方程，其通解为
\[
\]

（3）

将，得
\[
\]
\[
\]
将，阶跃响应
\[
\]
引用，即
\[
\]
$u_c$

根据，求得电容电流为
\[
\]
根据

$i_c$

线性电路，阶跃特性为
\[
\]
式中

现将响应为
\[
\]
再将，得到
\[
\]
$u_c'$所示

上述，表达式为
\[
\]
当

\section{一阶电路的冲激响应}

\section{正弦电源作用下的一阶电路}

\section{一阶电路的全响应}

\section{一阶电路暂态响应的一般形式}

\section{求解暂态响应的卷积积分法}

\section{二阶电路的暂态过程}

\section{状态变量分析}